\section{Discussion}

\subsection*{A reproducible framework rather than a single dataset finding}

The primary contribution of this work is an open-source, modular
framework for applying topological data analysis to longitudinal
scRNA-seq data, with integrated permutation null testing, cell-cycle
ablation, Wasserstein diagram comparison, and Mapper-based gene
attribution. The five datasets analyzed here are a deliberately
heterogeneous benchmark (two drugs, three biological scales, two PDX
models, one patient atlas, one 14-patient longitudinal cohort) chosen to stress-test reproducibility rather
than to maximize statistical significance on any single cohort. This
reframes the question from ``does drug tolerance produce topological
cyclicity?'' to ``is topological cyclicity a reproducible, quantitative
phenotype across biological systems that can be detected and compared
with a standardized pipeline?'' Our answer is yes: max $H_1$ increases
monotonically with osimertinib exposure in three independent systems
(PC9 cell line, YU-003 PDX, YU-006 PDX), and all six 1{,}000-cell
permutation tests in \cref{tab:validation} reach $p < 0.05$ (5/6 at
$p < 0.01$; the lone exception is the PC9 osimertinib baseline D0 at
$p = 0.030$).

\subsection*{Convergent topology, divergent gene programs}

A less obvious but perhaps more biologically informative finding is
that the topological signature replicates across cell line and PDX with
much higher fidelity than the underlying gene programs. Cell line
top-connected genes are dominated by canonical EMT markers (\emph{TPM1}, \emph{KRT8},
\emph{KRT19}), while PDX top-connected genes are a broader
mesenchymal/secretory signature (\emph{S100A4}, \emph{LGALS3}, \emph{COL1A1}, \emph{AGR2}). The
11-gene overlap between the two top-100 lists, although highly
significant statistically (hypergeometric $p = 1.25 \times 10^{-11}$),
is a minority of either list. We interpret this as evidence that the
cell-state plasticity underlying DTP emergence is realized through
partially different but related gene expression repertoires depending
on microenvironment --- an observation consistent with recent findings
on context-dependent EMT programs \citep{dongre2019}. The topological
statistic abstracts over these specific gene choices, making it a more
stable cross-system phenotype than any individual gene signature.

\subsection*{Relationship to existing literature}

Our work is complementary to and independent of \citet{oren2021}, who
used Watermelon lineage tracing to establish that PC9 DTPs contain
cycling subpopulations, and \citet{aissa2021}, who identified discrete
DTP transcriptional states in the source GSE134839 dataset. Our
framework provides an orthogonal method of detecting cyclic behavior
directly from the scRNA-seq transcriptome without lineage-tracing
infrastructure, and extends the signal to PDX data Oren et al.\ did not
analyze. Relative to \citet{rizvi2017} scTDA, which pioneered TDA on
mouse embryonic stem cell differentiation, we extend the framework to a
perturbation-response setting (rather than developmental) and introduce
controls (cell-cycle ablation, formal permutation null) absent from the
original formulation.

The broader non-genetic drug-tolerance literature --- YAP-mediated
senescent DTPs \citep{kurppa2020}, melanoma minimal-residual-disease
four-state plasticity \citep{rambow2018}, and the Marine et al.\
review of non-genetic resistance \citep{marine2020,shen2020} ---
provides the biological context. Our contribution against this
backdrop is a quantitative, reproducible statistic
(topologically invariant in loop existence, metric-dependent in
numerical scale; Supplementary Methods~S4) for detecting and
comparing such structures.

\subsection*{Limitations}

\paragraph{Platform and batch confounding.}
Cross-system comparisons mix Drop-seq (GSE134839) with 10x Chromium
(others); the cross-system ordering (cell line $<$ PDX $<$ patient)
is reported only qualitatively. Within-system comparisons are
platform-matched and survive Harmony batch correction
(\cref{fig:harmony}), ruling out within-dataset batch structure as
the source. A full scVI cross-platform integration
\citep{lopez2018} is deferred.

\paragraph{Permutation null specification.}
The gene-label null tests only for non-random point-cloud structure;
the complementary timepoint-label null on osi D14 vs.\ D0
($p < 0.02$; Results) directly tests drug dependence.

\paragraph{Patient-cohort validation.}
The three Maynard matched pre/post pairs (sign-test $p = 0.125$;
Results) are hypothesis-generating, not confirmatory; the pooled
14-patient TN$\to$PER$\to$PD monotonic increase is supportive but
observational. A prospective paired-biopsy study would be needed to
power the within-patient comparison.

\paragraph{Cell-cycle ablation stringency.}
Tirosh ablation removes 39/97 cell-cycle genes (standard but not
exhaustive --- quiescence and SASP markers are not in the Tirosh
list); the more demanding S/G2M-score regression (ratio 1.07--4.05
across five cohorts) and the PC1-filter Mapper (cell-cycle genes
co-rank with EMT markers) together indicate the topology is
multi-program rather than purely EMT, and is not a cell-cycle
artefact.

\subsection*{Implications and outlook}

If drug tolerance operates through topologically cyclic plasticity,
state-trapping combinations that collapse the cycle may outperform
single-state-targeting strategies. The cross-scale reproducibility of
the signature argues that max $H_1$ is a quantitative phenotype for
cell-state plasticity, suitable for benchmarking drug responses
across labs and platforms and likely to generalise to adjacent
perturbation-response settings.

\paragraph{Why might cyclic plasticity be adaptive?}
We speculate that a robust 1-cycle corresponds to a Waddington-style
limit cycle around a shallow basin \citep{huang2012,pisco2015},
making cyclic plasticity adaptive under fluctuating selection as
bet-hedging \citep{brock2009,gupta2011}; in this view osimertinib
both selects pre-existing cycles and induces new ones, although
distinguishing the two would require lineage tracing or a
controlled-clonal-bottleneck design. Adaptive therapy
\citep{gatenby2009,enriquez-navas2016} and state-trapping
combinations become testable hypotheses against this readout.
